% =====================================================================
%  2bat-ud4-14.tex  —  SESSIÓ 14: Estacions d'optimització
%  (sense preàmbul: s'inclou des de main.tex)
% =====================================================================
\horatitol{Sessió 14 — Estacions d'optimització}%
{Repàs actiu per quatre estacions, una per bloc de la unitat, i mapa final que connecta derivada, creixement, extrems i optimització.}

\subsection{Idea clau}

Repassem \textbf{tota la unitat} de manera activa, rotant per \textbf{quatre
estacions}, una per cada bloc treballat. Al final construirem un \textbf{mapa de la
unitat} que connecta les quatre idees —\textbf{derivada}, \textbf{creixement},
\textbf{màxims i mínims}, \textbf{optimització}— amb la lectura d'un model.

\begin{idea}
\textbf{Les quatre estacions}
\begin{enumerate}[leftmargin=1.6em, itemsep=2pt]
  \item \textbf{Càlcul de derivades:} derivar polinomis i avaluar la derivada en un punt.
  \item \textbf{Creixement i signe:} estudiar on una funció creix i decreix amb el signe
        de $f'$.
  \item \textbf{Màxims i mínims:} trobar i classificar els punts crítics.
  \item \textbf{Optimització:} plantejar i resoldre un problema de context social.
\end{enumerate}
\end{idea}

\subsection{Les estacions}

\begin{exemple}[Estació 1 — càlcul de derivades]
Per a $f(x)=x^3-4x^2+x-7$, la derivada és $f'(x)=3x^2-8x+1$ (terme a terme; el $-7$ es
deriva a $0$). Avaluant en $x=2$: $f'(2)=12-16+1=-3$.
\end{exemple}

\begin{exemple}[Estació 3 — màxims i mínims]
Per a $f(x)=x^3-12x+1$, $f'(x)=3x^2-12=3(x-2)(x+2)$, que s'anul·la a $x=-2$ i $x=2$. En
$x=-2$, $f'$ passa de $+$ a $-$: \textbf{màxim} $(-2,17)$. En $x=2$, de $-$ a $+$:
\textbf{mínim} $(2,-15)$.
\end{exemple}

\subsection{Exercicis resolts}

\begin{exresolt}[model de l'Estació 2: creixement i signe]
Estudia on creix i on decreix $f(x)=x^3-3x^2-9x$.
\solucio
$f'(x)=3x^2-6x-9=3(x-3)(x+1)$, s'anul·la a $x=-1$ i $x=3$. Signes: $(-\infty,-1)$
$f'(-2)=15>0$ creix; $(-1,3)$ $f'(0)=-9<0$ decreix; $(3,\infty)$ $f'(4)=15>0$ creix.
\resposta{Creix a $(-\infty,-1)$ i $(3,\infty)$; decreix a $(-1,3)$.}
\end{exresolt}

\begin{exresolt}[model de l'Estació 4: optimització]
Una entitat tanca un recinte rectangular contra una paret amb $20$ m de tanca (tres
costats). Si els costats perpendiculars a la paret fan $x$, troba les mides de màxima
superfície.
\solucio
El costat paral·lel a la paret fa $20-2x$, així que l'àrea és $A(x)=x(20-2x)=20x-2x^2$.
$A'(x)=20-4x=0\Rightarrow x=5$ (passa de $+$ a $-$: màxim). Mides: $5$ m i
$20-10=10$ m; superfície màxima $A(5)=50\ \text{m}^2$.
\resposta{$x=5$ m i $10$ m; superfície màxima $50\ \text{m}^2$.}
\end{exresolt}

\subsection{Exercicis per fer}

\noindent\textit{Una tasca per estació. Rota i resol-les totes.}

\begin{exercicis}
\begin{exlist}
  \item \textbf{(Estació 1)} Deriva $f(x)=2x^3-x$ i avalua $f'(1)$.
  \item \textbf{(Estació 1)} Deriva $g(x)=3x^2-5x+2$ i troba on $g'(x)=0$.
  \item \textbf{(Estació 2)} Estudia on creix i on decreix $f(x)=-x^2+10x$ i digues on
        té el punt més alt.
  \item \textbf{(Estació 3)} Troba i classifica els punts crítics de $f(x)=x^3-27x$.
  \item \textbf{(Estació 4)} Els ingressos d'una activitat (en euros) segons el preu $x$
        són $I(x)=-x^2+30x$. Troba el preu òptim i l'ingrés màxim.
  \item \textbf{(Mapa de la unitat)} Completa el mapa connectant cada concepte amb la
        seva idea:
        \begin{center}
        \renewcommand{\arraystretch}{1.4}
        \begin{tabular}{p{3.8cm} p{8.6cm}}
        \toprule
        \textbf{Concepte} & \textbf{Idea / lectura} \\
        \midrule
        Derivada $f'(x)$ & \rule{7.8cm}{0.4pt} \\
        Signe de $f'$ & \rule{7.8cm}{0.4pt} \\
        $f'(x)=0$ (punt crític) & \rule{7.8cm}{0.4pt} \\
        Optimització & \rule{7.8cm}{0.4pt} \\
        \bottomrule
        \end{tabular}
        \end{center}
\end{exlist}
\end{exercicis}

% ---------------------------------------------------------------------
\fullsolucions{Sessió 14}
\begin{solucions}
\begin{sollist}
  \item $f'(x)=6x^2-1$; $f'(1)=6-1=5$.
  \item $g'(x)=6x-5$; $g'(x)=0\Rightarrow x=\dfrac{5}{6}$.
  \item $f'(x)=-2x+10=0\Rightarrow x=5$. Creix a $(-\infty,5)$, decreix a $(5,\infty)$;
        punt més alt a $x=5$ ($f(5)=25$).
  \item $f'(x)=3x^2-27=3(x-3)(x+3)$, crítics $x=-3$ i $x=3$. Màxim a $(-3,54)$ ($f'$ de
        $+$ a $-$); mínim a $(3,-54)$ ($f'$ de $-$ a $+$).
  \item $I'(x)=-2x+30=0\Rightarrow x=15$ (màxim). Preu òptim $15\eur$; ingrés màxim
        $I(15)=-225+450=225\eur$.
  \item Resposta oberta. Idees: \textbf{derivada} = ritme de canvi / pendent a cada punt;
        \textbf{signe de $f'$} = si la funció creix ($+$) o decreix ($-$); \textbf{punt
        crític} ($f'=0$) = candidat a màxim o mínim (canvi de tendència);
        \textbf{optimització} = construir una funció i trobar-ne l'extrem per resoldre un
        problema real.
\end{sollist}
\end{solucions}
